\chapter{PENELITIAN TERKAIT}

\section{State of the Art Teknologi Pemantauan dan Pemasaran Ternak Berbasis IoT}

Pemanfaatan Internet of Things (IoT) pada sektor peternakan terus berkembang dari pencatatan manual menuju pengumpulan data otomatis yang terhubung dengan sistem informasi. Farooq dkk. (2022) memetakan peran IoT dalam pertanian untuk mewujudkan lingkungan peternakan cerdas (\textit{smart livestock}). Pada tingkat implementasi, sejumlah penelitian menjadikan bobot badan sebagai parameter utama pemantauan (Kushartadi dkk., 2023; Efendi dkk., 2025; Sumiharto dan Malik, 2026). Perubahan bobot dari waktu ke waktu menggambarkan laju pertumbuhan hewan, sehingga data yang tercatat secara berkala dapat menjadi dasar pengambilan keputusan pemeliharaan (Efendi dkk., 2025).

Kushartadi dkk. (2023) merancang sistem pemantauan ternak yang memadukan pengukur bobot cerdas dengan penanda Bluetooth Low Energy (BLE) yang dikenakan pada hewan sebagai identitas. Data bobot dan kesehatan dikirim ke sistem awan. Sistem tersebut dibangun dengan mikrokontroler ESP32, sensor beban (\textit{load cell}), BLE, dan protokol Message Queuing Telemetry Transport (MQTT), sehingga pemantauan dapat dilakukan secara waktu nyata dan jarak jauh. Susunan perangkat keras serupa digunakan oleh Sumiharto dan Malik (2026) di Universitas Gadjah Mada, yaitu perangkat berbasis ESP32 yang membaca identitas ternak melalui RFID dan bobot melalui \textit{load cell}. Data dikirim dalam format JSON ke peladen untuk divalidasi sebelum disimpan pada basis data PostgreSQL yang dikelola melalui Supabase.

Penelitian Sumiharto dan Malik (2026) juga menunjukkan bahwa keandalan sistem tidak hanya ditentukan oleh perangkat keras. Mereka menerapkan aplikasi \textit{full-stack} Next.js dengan Supabase Auth dan kebijakan Row Level Security (RLS) untuk membatasi akses data berdasarkan kepemilikan pengguna. Setiap data yang masuk melalui verifikasi kunci API, validasi struktur data, dan pemeriksaan status perangkat. Pengujian beban pada rentang 40 hingga 95 permintaan per detik menunjukkan bahwa sistem stabil hingga 90 permintaan per detik tanpa temuan kerentanan tingkat tinggi. Dari sisi pemanfaatan data, Efendi dkk. (2025) mengumpulkan bobot mingguan 15 ekor domba selama enam bulan menggunakan Arduino dan ESP32. Pertumbuhan setiap hewan kemudian dimodelkan dengan regresi linear sederhana, yang menghasilkan nilai $R^2$ di atas 99\% pada sebagian besar ternak. Hasil tersebut disajikan melalui antarmuka telepon pintar.

Di Indonesia, sistem serupa telah diterapkan pada ternak kambing. Amrulloh dkk. (2024) mengembangkan timbangan kambing berbasis IoT yang terdiri atas mikrokontroler Arduino Uno, sensor berat, dan sensor RFID untuk mengirimkan identitas serta bobot ternak ke basis data melalui API RESTful. Data tersebut ditampilkan pada aplikasi web berbasis Laravel dan aplikasi seluler berbasis React Native untuk memantau perkembangan ternak di Desa Dermaji, Kabupaten Banyumas. Menurut penelitian tersebut, data bobot dikirim ketika angka pada monitor timbangan tidak lagi berubah secara drastis, yang dianggap menandakan posisi kambing sudah stabil. Sholicha dkk. (2023) merancang aplikasi Smart Goat untuk pencatatan ternak pada program penggemukan kambing. Aplikasi tersebut dilengkapi sensor RFID, timbangan digital, serta pemantauan suhu, kelembapan, intensitas cahaya, dan gas amonia pada kandang.

Kebutuhan terhadap alat semacam ini muncul karena penimbangan langsung sering dianggap merepotkan. Amrulloh dkk. (2024), dengan merujuk penelitian sebelumnya, menyebutkan bahwa pengukuran bobot kambing di sebagian peternakan masih dilakukan dengan mengangkat kambing ke atas timbangan gantung menggunakan karung. Sebagai alternatif, penduga bobot berbasis citra juga dikembangkan. Shuvo dkk. (2026) membangun aplikasi web progresif (\textit{Progressive Web App}, PWA) yang dapat berjalan secara luring dan menduga bobot sapi dari foto serta ukuran tubuh. Penelitian tersebut melaporkan galat absolut rata-rata (\textit{Mean Absolute Error}, MAE) sebesar 7,52 kg pada lima ekor sapi uji dengan bobot 213--412 kg. Timbangan elektronik bersertifikat dengan ketelitian $\pm 0{,}5$ kg tetap digunakan sebagai acuan. Kulesz (2026) juga melaporkan PWA untuk peternak kambing dengan fitur penilaian bobot berbasis foto yang berjalan sepenuhnya luring. Temuan ini menunjukkan bahwa timbangan tetap diperlukan sebagai acuan pengukuran langsung.

Data bobot yang akurat memiliki nilai ekonomi langsung dalam transaksi. Rianto dkk. (2026) menelaah 90 ekor kambing kurban di Kabupaten Semarang dan menemukan korelasi sangat kuat antara harga per ekor dan bobot badan ($r = 0{,}848$). Bobot menjelaskan 71,89\% variasi harga, sedangkan setiap kenaikan 1 kg bobot berkaitan dengan kenaikan harga sekitar Rp79.565 per ekor pada sampel tersebut. Angka ini merupakan hasil pada sampel penelitian dan tidak boleh digeneralisasi ke seluruh pasar tanpa pengujian tambahan. Meskipun demikian, penimbangan tidak selalu dilakukan. Rianto dkk. (2026) menduga bobot dari lingkar dada dan panjang badan pada hewan yang tidak dapat ditimbang. Septyan dkk. (2023) menyebutkan bahwa penjualan sapi potong di Indonesia umumnya masih mengandalkan taksiran penampilan badan. Pada sapi Bali, rumus pendugaan Djagra dan Lambourne yang telah diuji masih menyimpang sekitar 8,5\% hingga 8,7\% dari hasil timbangan digital.

Dari sisi pemasaran, platform digital untuk jual beli kambing telah dikembangkan di Indonesia. Waluyo dkk. (2018) membangun aplikasi \textit{m-commerce} berbasis Android sebagai media informasi dan penjualan kambing yang mempertemukan pembeli dengan peternak. Penelitian tersebut dilatarbelakangi oleh keterbatasan waktu operasional pasar hewan, peran perantara, dan biaya operasional yang dapat meningkatkan harga di tingkat pembeli. Oktaviani dkk. (2023) merancang \textit{e-marketplace} bagi kelompok ternak kambing yang melayani jual beli, lelang, dan bagi hasil ternak. Penelitian tersebut berangkat dari harga jual hasil ternak yang berada di bawah harga pasar. Dengan demikian, teknologi pengukuran otomatis, pengelolaan data berbasis awan, dan platform pemasaran digital telah berkembang, tetapi umumnya dikembangkan secara terpisah. Keterkaitan ketiganya dan celah penelitian diuraikan pada subbab berikutnya.

\section{Tinjauan Penelitian Terdahulu}

Penelitian terdahulu yang berkaitan dengan topik ini dapat dikelompokkan ke dalam tiga kelompok. Kelompok pertama adalah sistem timbangan ternak berbasis IoT yang terhubung dengan platform pemantauan. Kelompok kedua adalah aplikasi pencatatan ternak yang memanfaatkan timbangan sebagai salah satu sumber data. Kelompok ketiga adalah penduga bobot ternak tanpa timbangan, misalnya berbasis citra. Pemahaman mengenai capaian dan batasan masing-masing kelompok diperlukan untuk menentukan posisi penelitian ini. Perbandingan penelitian terdahulu dengan penelitian yang diusulkan dirangkum pada Tabel~\ref{tab:penelitian-terdahulu}.

\begin{table}[H]
  \centering
  \caption{Perbandingan penelitian terdahulu dan penelitian yang diusulkan}
  \label{tab:penelitian-terdahulu}
  \small
  \begin{tabularx}{\textwidth}{p{0.20\textwidth} p{0.23\textwidth} X X}
    \toprule
    \textbf{Peneliti (Tahun)} & \textbf{Fokus penelitian} & \textbf{Keunggulan sistem} & \textbf{Batasan atau kekurangan} \\
    \midrule
    Amrulloh dkk. (2024) & Timbangan kambing IoT (Arduino Uno, sensor berat, RFID) dengan aplikasi web Laravel dan aplikasi seluler React Native. & Tiga platform terintegrasi melalui API RESTful; data bobot dikirim saat angka dianggap stabil. & Pengujian terbatas pada fungsionalitas \textit{black-box}; akurasi terhadap timbangan acuan belum diuji; kriteria stabil bersifat kualitatif. \\
    Sholicha dkk. (2023) & Aplikasi Smart Goat untuk pencatatan ternak pada program penggemukan kambing. & Manajemen data ternak, pemantauan bobot, dan pemantauan kondisi kandang dengan RFID serta timbangan digital. & Berfokus pada pencatatan dan pemantauan kandang; evaluasi akurasi terhadap timbangan acuan belum dilaporkan. \\
    Kushartadi dkk. (2023) & Pemantauan ternak waktu nyata dengan pengukur bobot cerdas, BLE, ESP32, \textit{load cell}, dan MQTT. & Data bobot dan kesehatan dikirim ke awan serta dapat dipantau jarak jauh. & Tidak membahas penanganan hewan yang bergerak saat penimbangan. \\
    Sumiharto dan Malik (2026) & Sistem pemantauan bobot ternak berbasis web dengan timbangan IoT, Next.js, dan Supabase. & Arsitektur web--awan--IoT dengan kontrol akses RLS; stabil hingga 90 permintaan per detik pada pengujian yang dilaporkan. & Pengujian menitikberatkan kinerja dan keamanan sistem; sistem berbasis web. \\
    Rahayu dkk. (2023) & Sistem penimbangan ternak berbasis IoT untuk marketplace Caribi pada skema investasi ternak Maro. & Membaca RFID, menimbang ternak, dan mengirim bobot ke peladen. & Berorientasi pada skema investasi bagi hasil, bukan kualitas data penimbangan. \\
    Shuvo dkk. (2026) & PWA luring untuk menduga bobot sapi dari citra dan ukuran tubuh. & Berjalan tanpa internet; menyediakan skor kepercayaan dan uji kebergunaan. & Berupa pendugaan, bukan pengukuran langsung; objeknya sapi pada fasilitas dengan jumlah hewan uji terbatas. \\
    Penelitian ini & Timbangan IoT berpalang pintu (\textit{gated}) dan aplikasi untuk pengumpulan data kambing. & Deteksi bobot stabil yang divalidasi terhadap timbangan acuan, skor kualitas tiap penimbangan, serta jalur data \textit{offline-first} yang diuji. & Diisi sesuai batasan produk pada Bab I. \\
    \bottomrule
  \end{tabularx}
\end{table}

Berdasarkan Tabel~\ref{tab:penelitian-terdahulu}, kelompok sistem timbangan IoT telah menunjukkan bahwa bobot ternak dapat dikumpulkan secara otomatis dan dikirim ke aplikasi. Meskipun demikian, pengujian pada penelitian tersebut cenderung diarahkan pada fungsi dan kinerja sistem. Amrulloh dkk. (2024), misalnya, menguji fungsionalitas dengan pengujian \textit{black-box} dan menetapkan kondisi stabil berdasarkan perubahan angka pada monitor tanpa kriteria kuantitatif yang dapat direplikasi. Sumiharto dan Malik (2026) menguji beban dan keamanan sistem. Keduanya belum menguji seberapa dekat bobot yang dikirim timbangan dengan timbangan acuan pada hewan hidup yang bergerak.

Penelitian Smart Goat (Sholicha dkk., 2023) dan sistem pada platform investasi (Amrulloh dkk., 2024; Rahayu dkk., 2023) menempatkan timbangan sebagai salah satu sumber data bagi aplikasi yang tujuan utamanya adalah pencatatan atau investasi. Kualitas pembacaan timbangan itu sendiri belum menjadi objek utama penelitian.

Kelompok ketiga menggunakan prinsip pengukuran yang berbeda. Shuvo dkk. (2026) dan Kulesz (2026) menunjukkan bahwa PWA luring dapat dijalankan di lapangan untuk pendugaan bobot dari foto. Akan tetapi, pendugaan tersebut diuji pada objek dan jumlah hewan yang terbatas, dan penelitian Shuvo dkk. (2026) tetap membandingkan hasilnya dengan timbangan elektronik. Hal ini memperkuat posisi timbangan sebagai acuan yang masih perlu dibuat andal pada kambing. Perbedaan tersebut membentuk celah penelitian yang diuraikan pada subbab berikutnya.

\section{Posisi dan Kesenjangan Penelitian (\textit{Research Gap})}

Tinjauan literatur menunjukkan bahwa sistem timbangan ternak berbasis IoT dengan aplikasi seluler dan web telah dibangun, termasuk untuk kambing. Berdasarkan penelusuran pustaka yang digunakan dalam proposal ini, kualitas pengukuran pada sistem-sistem tersebut belum menjadi pokok penelitian. Kondisi bobot stabil ditentukan secara kualitatif, belum ada penilaian kualitas untuk setiap pembacaan, dan perilaku jalur data pada kondisi sinyal buruk belum diuji. Padahal, bobot badan berkaitan dengan nilai ternak (Rianto dkk., 2026), sedangkan pendugaan tanpa timbangan masih memiliki penyimpangan sekitar 8,5\% hingga 8,7\% pada rumus yang diuji pada sapi Bali (Septyan dkk., 2023).

Penelitian ini diposisikan untuk mengisi celah tersebut melalui tiga kesenjangan. Kesenjangan pertama adalah metode deteksi bobot stabil pada timbangan berpalang pintu untuk kambing. Sinyal \textit{load cell} dibaca secara kontinu selama kambing berada pada timbangan, kemudian jendela pembacaan dengan variansi terendah dipilih dan nilai mediannya dilaporkan sebagai bobot. Galat, bias, dan keterulangan metode ini diukur terhadap timbangan acuan pada kambing hidup, sehingga hasilnya dapat dibandingkan dengan galat pendugaan berbasis ukuran tubuh.

Kesenjangan kedua adalah lapisan kualitas data pada setiap penimbangan. Setiap pembacaan diberi skor kualitas berdasarkan variansi sinyal dan lamanya kambing diam. Aplikasi meminta penimbangan ulang apabila skor rendah. Gagasan ini sejalan dengan skor kepercayaan pada Smart-Herbo yang menandai estimasi untuk ditinjau ulang (Shuvo dkk., 2026), tetapi diterapkan pada pembacaan timbangan kambing, bukan pada estimasi citra sapi. Pada riwayat bobot setiap ekor, pencilan juga dapat disaring. Algoritma Kalman filter with impulse noisy outliers (kfino) telah divalidasi pada data bobot hasil \textit{walk-over-weighing} yang banyak memuat pencilan (Gonzalez-Garcia dkk., 2025). Dalam penelitian ini, kfino dibandingkan dengan filter median sederhana.

Kesenjangan ketiga adalah jalur data \textit{offline-first} yang diuji. Kandang sering memiliki sinyal lemah, sehingga aplikasi menyimpan pembacaan secara lokal dan menyinkronkannya setelah koneksi pulih. Aplikasi luring pada penelitian terdahulu, yaitu Smart-Herbo dan DigiZiege, bekerja dengan pendugaan berbasis citra, bukan dengan timbangan (Shuvo dkk., 2026; Kulesz, 2026). Dalam penelitian ini, jalur data dari timbangan hingga basis data diuji melalui skenario pemutusan koneksi dengan indikator tingkat keberhasilan sinkronisasi, kehilangan data, dan duplikasi data.

Dengan demikian, penelitian ini tidak bertumpu pada kebaruan teknologi pembangun aplikasi, melainkan pada kualitas dan keandalan data penimbangan kambing yang diukur secara terukur. Data yang akurat dan terjamin kualitasnya dapat menjadi dasar pengembangan lanjutan, misalnya platform jual beli ternak yang menampilkan data terukur kepada calon pembeli.
